\subsection{\texorpdfstring{The relation \(\sum_{i} e_i f_i \leqslant n\)}{The relation Σ ei fi ≤ n}}

Let K be a field, v a valuation on K and L a finite extension of K of degree n.~Let \((v_{1}^{\prime}, \ldots, v_{r}^{\prime})\) be valuations on L extending v no two of which are equivalent; if they are independent (which is always the case if v is of height 1), then

\[
\sum_ {i = 1} ^ {r} e (v _ {i} ^ {\prime} / v) f (v _ {i} ^ {\prime} / v) \leqslant n
\]

by Proposition 2 (formulae (6) and (7)). We shall see that this result is true in the general case. To be precise:

THEOREM 1. Let K be a field, v a valuation on K and L a finite extension of K of degree n.~Then:

\begin{enumerate}
\def\labelenumi{(\alph{enumi})}
\item
  Every complete system \((v_i')_{i\in I}\) of extensions of \(v\) to \(\mathbf{L}\) is finite.
\item
  \(\sum_{i\in I}e(v_i^{\prime}|v)f(v_i^{\prime}|v)\leqslant n\) and a fortiori Card(I) \(\leqslant n.\)
\item
  No two of the rings of the \(v_{i}'\) are comparable with respect to inclusion.
\end{enumerate}

Since the theorem is trivial if \(v\) is improper, we shall assume that \(v\) is not improper. Let \((v_1', \ldots, v_s')\) be any finite family of valuations on \(L\) extending \(v\) , no two of which are equivalent. We shall first prove that \(\sum_{i=1}^{s} e(v_i'/v)f(v_i'/v) \leqslant n\) . This will prove (a) and (b).

We argue by induction on s and suppose therefore that the inequality has been established for the case of 0, 1, \ldots, s - 1 valuations. We distinguish two cases.

\begin{enumerate}
\def\labelenumi{(\arabic{enumi})}
\tightlist
\item
  Suppose that there exist at least two independent valuations \(v_i'\) . Then there exists (§ 7, no. 2, Remark 1) a partition \([1, s] = I_1 \cup \cdots \cup I_t\) of \([1, s]\) such that:
\end{enumerate}

\begin{enumerate}
\def\labelenumi{(\roman{enumi})}
\item
  for \(v_{i}^{\prime}\) and \(v_{j}^{\prime}\) to be dependent, it is necessary and sufficient that \(i\) and \(j\) belong to the same \(I_{k}\) ;
\item
  \(\operatorname{Card}(\mathbf{I}_k) < s\) for all \(k\) .
\end{enumerate}

We choose in each \(I_{k}\) an index \(i(k)\) . Let \(\hat{L}_{i(k)}\) denote the completion of L with respect to \(v_{i(k)}^{\prime}\) and \(n(k) = [\hat{L}_{i(k)} : \hat{K}]\) . For all \(i \in I_{k}\) , \(v_{i}^{\prime}\) defines on L the same topology as \(v_{i(k)}^{\prime}\) (§ 7, no. 2, Proposition 3) and hence may be extended to a valuation \(\hat{v}_{i}^{\prime}\) on \(\hat{L}_{i(k)}\) whose restriction to \(\hat{K}\) is \(\hat{v}\) . Since no two of the \(v_{i}^{\prime}\) for \(i \in I_{k}\) are equivalent, the same is true of the \(\hat{v}_{i}^{\prime}\) . The induction hypothesis applied to the ordered pair ( \(\hat{K}, \hat{L}_{i(k)}\) ) shows, by virtue of Proposition 2 (a), formula (4), that \(\sum_{i\in I_k}e(v_i^{\prime}|v)f(v_i^{\prime}|v)\leqslant n(k)\) . As \(\sum_{k = 1}^{t}n(k)\leqslant n\) (Proposition 2 (b), formula

(7)), certainly \(\sum_{i=1}^{s} e(v_i'/v) f(v_i'/v) \leqslant n\) .

\begin{enumerate}
\def\labelenumi{(\arabic{enumi})}
\setcounter{enumi}{1}
\tightlist
\item
  We now pass to the case where any two of the \(v_i'\) are dependent. Let \(A_i'\) be the ring of \(v_i' (1 \leqslant i \leqslant s)\) ; writing A for the ring of \(v, A_i' \cap K = A\) for all \(i\) . Let \(B'\) be the subring of L generated by \(A_1', \ldots, A_s'\) ; we write \(B = B' \cap K\) ; then \(B \supset A\) . Then B is the ring of a valuation \(w\) on K and \(B'\) the ring of a valuation \(w'\) which is not improper and extends \(w\) ( \(\S 7\) , no. 2, Proposition 4); the field \(\kappa(B')\) is an extension of \(\kappa(B)\) of degree \(f(w'/w)\) . Consider the canonical images \(\bar{A}_i', \bar{A}\) of \(A_i'\) and A in \(\kappa(B')\) ; then \(\bar{A}\) is the ring of a valuation \(\bar{v}\) on \(\kappa(B)\) and the \(\bar{A}_i'\) are rings of valuations \(\bar{v}_i'\) on \(\kappa(B')\) extending \(\bar{v}\) . As the \(A_i'\) generate \(B'\) , the \(\bar{A}_i'\) generate \(\kappa(B')\) and hence the \(\bar{v}_i'\) are not all dependent ( \(\S 7\) , no. 2, Proposition 4). From the first part of the proof,
\end{enumerate}

\[
\sum_ {i = 1} ^ {s} e \left(\bar {v} _ {i} ^ {\prime} / \bar {v}\right) f \left(\bar {v} _ {i} ^ {\prime} / \bar {v}\right) \leqslant \left[ \kappa \left(\mathbf {B} ^ {\prime}\right): \kappa (\mathbf {B}) \right] = f \left(w ^ {\prime} / w\right)
\]

and hence

\[
\sum_ {i = 1} ^ {s} e (w ^ {\prime} / w) e (\bar {v} _ {i} ^ {\prime} / \bar {v}) f (\bar {v} _ {i} ^ {\prime} / \bar {v}) \leqslant e (w ^ {\prime} / w) f (w ^ {\prime} / w) \leqslant n \quad \text {(no. 1, Lemma 1)}.
\]

The proof of (a) and (b) will therefore be completed if we prove that

\[
f (\bar {v} _ {i} ^ {\prime} / \bar {v}) = f (\bar {v} _ {i} ^ {\prime} / \bar {v}), \quad e (w ^ {\prime} / w) e (\bar {v} _ {i} ^ {\prime} / \bar {v}) = e (v _ {i} ^ {\prime} / v).\tag{8}
\]

For this, we note that v and \(\bar{v}\) (resp. \(v_{i}^{\prime}\) and \(\bar{v}_{i}^{\prime}\) ) have the same residue field (§ 4, no. 1, Corollary to Proposition 2); this proves the first equation. For the second there is, by virtue of the Remark in § 4, no. 3, the following commutative diagram, where the rows are exact sequences and the vertical arrows represent canonical injections:

\[
\begin{array}{c} 0 \to \Gamma_ {\bar {v}} \to \Gamma_ {v} \to \Gamma_ {w} \to 0 \\ \downarrow \quad \downarrow \quad \downarrow \\ 0 \to \Gamma_ {\bar {v} _ {i} ^ {\prime}} \to \Gamma_ {v _ {i} ^ {\prime}} \to \Gamma_ {w ^ {\prime}} \to 0 \end{array}
\]

We deduce that there is an exact sequence

\[
0 \rightarrow \Gamma_ {\bar {v} _ {i} ^ {\prime}} / \Gamma_ {\bar {v}} \rightarrow \Gamma_ {v _ {i} ^ {\prime}} / \Gamma_ {v} \rightarrow \Gamma_ {w ^ {\prime}} / \Gamma_ {w} \rightarrow 0
\]

by Chapter I, § 1, no. 4, Proposition 2, which proves the second formula (8). To complete the proof of Theorem 1, it remains to prove (c). If the ring of \(v_i'\) contains that of \(v_j'\) , \(\Gamma_{v_i'}\) is identified with a quotient group \(\Gamma_{v_j'} / H\) , \(H\) being an isolated subgroup (§ 4, no. 3). As the composite canonical mapping

\[
\Gamma_ {v} \rightarrow \Gamma_ {v _ {j} ^ {\prime}} \rightarrow \Gamma_ {v _ {j} ^ {\prime}} / H = \Gamma_ {v _ {i} ^ {\prime}}
\]

is injective, \(H \cap \Gamma_{v} = \{0\}\) , whence \(H = \{0\}\) (Lemma 3, no. 1). Then \(v_{i}'\) and \(v_{j}'\) are equivalent, whence i = j.

Remark. The intersection C of the rings \(A_{i}^{\prime}\) of the valuations \(v_{i}^{\prime}\) ( \(i \in I\) ) is the integral closure of A in L ( \(\S 1\) , no. 3, Corollary 3 to Theorem 3); it follows moreover from (c) and \(\S 7\) , no. 1, Propositions 1 and 2 that C is a semi-local ring, that its maximal ideals are the intersections \(m_{i} = C \cap m(A_{i}^{\prime})\) and that \(A_{i}^{\prime} = C_{m}\) for all \(i \in I\) .

